\documentclass[a4paper, 12pt]{article}  %set doc type and font size
\usepackage{fontspec}                   %allow setting fonts
\usepackage{xeCJK}         
\usepackage{amsmath}
\newtheorem{theorem}{Theorem}
\setCJKmainfont{Noto Sans CJK SC}             %call chinese package xeCJK
% \setmainfont{Times New Roman}           %set English font

% \XeTeXlinebreaklocale "zh"
% \XeTeXlinebreakskip = 0pt plus 1pt


\usepackage{algorithm}
\usepackage[noend]{algpseudocode}

\usepackage{mdframed}    % for frames
\usepackage{xcolor}      % for colors

\usepackage{amsthm}

\usepackage{graphicx}

\usepackage[most]{tcolorbox}
\tcbset{
  theorem style/.style={
    enhanced,
colback=blue!2!white,
colframe=blue!40!white,
    fonttitle=\bfseries,
    boxrule=1pt,
    arc=3pt,
    outer arc=3pt,
    coltitle=black,
  }
}

% 定義「帶框定理」環境
\newtcbtheorem[number within=section]{mytheorem}{Theorem}%
{theorem style=blue}{th}


\newtcbtheorem[number within=section]{problem}{Problem}%
{enhanced,
 colback=blue!3!white,
 colframe=blue!50!white,
 fonttitle=\bfseries,
 boxrule=1pt,
 arc=3pt,
 coltitle=black}{prob}



\title{電腦對局理論 Homework 2}
\author{蔡平樂}
\date{\today}


\begin{document}
\maketitle

\section{Implementation detail}
The following is the design be used at the final submission of the MCTS agent.

\subsection{Simulation and improvement}
For each single MCTS iteration, the agent will perform $10$ simulations, and use a PieceValue table to evaluate each move.

\begin{algorithm}
  \begin{algorithmic}[1]
    \Procedure{MoveEvaluate}{move}
      \State$s\gets move.from$
      \State$e\gets move.to$
      \If{$e\neq Null$}
        \State$capture\_move\_value \gets PieceValue(e)$
        \If{$s\neq{\rm Cannon}$}
          \State$capture\_move\_value += \frac{PieceValue(s)}{2}$
        \EndIf
        \State \textbf{return $C\times capture\_move\_value$}
      \EndIf
      \State \textbf{return $basic\_move\_value$}
    \EndProcedure
  \end{algorithmic}
\end{algorithm}

$basic\_move\_value$ and $C$ are constants, which are 


\subsection{RAVE}
The Rapid Action Value Estimation method only being used at the leaf.
More precisely, I create two types of Nodes, Nodes and Nodes\_AMAF, then update them respectively at a simulation

\begin{algorithm}
    \begin{algorithmic}
        \Procedure{UpdateNode}{node, sim\_record}
            
        \EndProcedure
    \end{algorithmic}
\end{algorithm}

\subsection{Evaluation function design}

\section{Experiments}



\end{document}